% Generated by records_reorg.py assemble (records-reorg/); edit the data and re-run, not this file.
\subsection{The detector-response mismatch}\label{recsec:supplementary-response}

% Moved verbatim from app_response_mismatch.tex @ 84da419c8 in the compression under DECISION-20261120a (qualifier and procedure detail
% may move to the records with a pointer); move-manifest-je JE-19.
% [records reorg] former heading: \subsection{Detector-response mismatch (App.~C): detail moved from the note}
\paragraph{Detector-response mismatch (App.~C): detail moved from the note.}

% [records reorg] former heading: \subsubsection{Reconstruction-only test}
\paragraph{Reconstruction-only test.}
The nominal reconstructed simulation, all truth coordinates, event weights, and event membership stay fixed. A documented hadronic-response variation or an explicit quantitative assumption should justify the amplitude of the reconstruction-only shift $\epsilon$; the overall factor in the reconstructed-$\Eavail$ definition, Eq.~\eqref{eq:eavail-reco-implemented}, is not by itself a sufficient basis.
Outputs include the full five-dimensional unfolded/reference residual, one-dimensional $\Eavail$ and $W$ projections, and an $(\Eavail,W)$ ratio map. The implementation runs only as a pseudo-data closure and only on the five observables, and cannot be combined with truth-reweight injection or bootstrapping. It labels every output as a diagnostic that must not be quoted.

% [records reorg] former heading: \subsubsection{Follow-up and interpretation}
\paragraph{Follow-up and interpretation.}
The current per-event outputs hold only scalar observables, not the low-level event record that the follow-up would alter, so the follow-up requires dedicated diagnostic output from the event processing.
A small observed bias would support stability of the $(\Eavail,W)$ interpretation only within the tested response envelope. A large bias would require a decision on the response model or an uncertainty component, and would not be absorbed silently into a non-closure error.

\subsection{Figure and table detail for Sections 2, 3 and 5 and the appendices}\label{recsec:supplementary-captions}

% [records reorg] former heading: \subsection{Caption detail (caption pass, Parts 2, 3, 5 and appendices)}
\paragraph{Caption detail (caption pass, Parts 2, 3, 5 and appendices).}
% The director's ruling of 2026-11-20 (option (b)): secondary caption sentences outside Part 4, moved verbatim (rows
% CV-CAPR.* in history-cuts/move-manifest-cv-captions-20261120.tsv). Bracketed words are context, not the note's.
% From app_negweight.tex:87 (draft d55f7aad6)
[Caption of the note's Fig.~\ref{fig:negweight-ratio}, continued:] [the low-$p_T$, high-background edge bin of the negative-weight/purity ratio map] Background is largest in the edge bin, and the $\max(0,\cdot)$/structural floor of \S\ref{sec:negweight-estimator} governs the result there.
% From sec_experiment.tex:52 (draft d55f7aad6)
[Caption of the note's Fig.~\ref{fig:control-corner}, continued:] [among the simulation's pairwise correlations of the five reconstructed observables] The strongest correlations share recoil energy, $q_{3}$ with $\Eavail$ ($r\approx0.81$) and $W$ with $\Eavail$ ($r\approx0.73$), or the muon transverse kick, $q_{3}$ with $\pT$ ($r\approx0.59$).
% From sec_experiment.tex:52 (draft d55f7aad6)
[Caption of the note's Fig.~\ref{fig:control-corner}, continued:] Apart from $\Eavail$ with $\pT$ ($r\approx0.34$), $\Eavail$ and $W$ are nearly uncorrelated with the muon variables.
% From sec_joint.tex:85 (draft d55f7aad6)
[Caption of the note's the earlier note's table of the joint generator tests, continued:] [the six conditions under which every joint test is rejected] The conditions concern (i) the null truth below the fine grid, (ii) the detector model, (iii) the interaction-model prior, (iv) the simulation size, (v) the neutrino-energy range and (vi) the predictions' own finite simulated samples (\S\ref{sec:joint5d-conditions}).
% From sec_joint.tex:111 (draft d55f7aad6)
[Caption of the note's the earlier note's figure of the joint tests' null distributions, continued:] Sufficient inputs: App.~\ref{app:release}.
% From sec_results.tex:111 (draft d55f7aad6)
[Caption of the note's Fig.~\ref{fig:ptslices}, continued:] Only the panel format follows Fig.~13 of Ref.~\cite{MINERvA:2021owq}.
% From sec_results.tex:111 (draft d55f7aad6)
[Caption of the note's Fig.~\ref{fig:ptslices}, continued:] Published-data comparisons are in Figs.~\ref{fig:xsec} and~\ref{fig:pull} and the earlier note's comparison table.
% From sec_results.tex:193 (draft d55f7aad6)
[Caption of the note's the earlier note's figure of the model-ratio maps, continued:] [our result's and the published data's ratio maps to MINERvA Tune~v1] The maps are close because both results depart from the tune in the same places by similar amounts.
% From selection_features.tex:6 (draft d55f7aad6)
[Caption of the note's the earlier note's table of selection and classifier features, continued:] Extended fiducial domain: App.~\ref{sec:fps}.
